\section{Local equality and uniform frequency bounds over the complex field}
\label{sec:hecke-original-upper}

Let $K=\mathbb Q(\sqrt{-3})$, let $\chi$ be a fixed primitive finite
Hecke character of conductor $\mathfrak f$, and put
\[
 c=\frac{2\pi}{\sqrt{3N\mathfrak f}},\qquad
 (\mathcal H_cf)(u)=c\int_0^\infty
                   J_0(2c\sqrt{uv})f(v)\,dv.
\]
The raw Hilbert space is $L^2((0,\infty),du)$.
The transform $\mathcal H_c$ is a unitary involution. Write
\[
 (Uf)(t)=e^{t/2}f(e^t),\qquad H_D^s=U^{-1}H^s(\mathbb R),
 \qquad (D_rf)(u)=r^{-1/2}f(u/r).
\]
Our Fourier convention on logarithmic space is
$\widehat g(\tau)=\int g(t)e^{i\tau t}dt$, with Plancherel measure
$d\tau/(2\pi)$. Let $C_a$ be the multiplier with symbol $a$.
The theta map has the closed multiplier realization
\[
 \widehat{E_\chi f}(\tau)
       =L_K(1/2+i\tau,\chi)\widehat{Uf}(\tau),
 \qquad
 E_\chi\mathcal H_c=\varepsilon_\chi J E_{\bar\chi},
 \quad (Jg)(t)=g(-t),\quad |\varepsilon_\chi|=1.
\]

Put $L=\log\Lambda$, $P_\Lambda=1_{[0,\Lambda]}$,
$B_\Lambda=\mathcal H_cP_\Lambda\mathcal H_c$, and
\[
 \Pi_\Lambda=
 1_{[1-\Lambda^{-56},1]}
       (P_\Lambda\mathcal H_cP_\Lambda\mathcal H_cP_\Lambda).
\]
For the principal character use
$\operatorname{Ran}\Pi_\Lambda\cap\ker I\cap\ker\operatorname{ev}_0$,
where $If=\int_0^\infty f(u)du$. Let $Q_\Lambda$ be the orthogonal
projection onto the corresponding theta image, with its full Gram
normalization, and set
\[
 S_L=1_{[-L,L]},\qquad D_L=S_L-S_LQ_\Lambda S_L.
\]
If $z_i$ is any orthonormal basis of that image, define
\[
 d\mu_\Lambda(\tau)=
 \left(2L-\sum_i|\widehat{S_Lz_i}(\tau)|^2\right)
                                      \frac{d\tau}{2\pi}.
\]

\paragraph{The estimates already established.}
The preceding sections give the following fixed-character estimates.
\begin{enumerate}
\item Proposition~\ref{prop:hankel-high-cutoff} gives selected rank
$O(\Lambda)$, and Proposition~\ref{prop:hecke-crude-gram} gives
arithmetic image Gram inverse $O(\Lambda^2)$.
\item Theorem~\ref{thm:hecke-deep-leakage} bounds the normalized theta
column outside $[-L,L]$ in Hilbert--Schmidt norm by $O(\Lambda^{-17/8})$.
\item Every fixed finite packet of nontrivial zero jets is asymptotically
annihilated by $Q_\Lambda$. More precisely, the consequence used here is
$\|Q_\Lambda P_{F,L}\|=O_{F,N}(L^{-N})$ for every fixed packet $F$
and every $N$, where $P_{F,L}$ projects onto its window jets. At a zero
$\rho=1/2+\alpha+i\gamma$, the jet generators are
\[
 \begin{cases}
 1_{[-L,L]}e^{(\alpha-i\gamma)(t-L)}(t-L)^j,&\alpha>0,\\
 1_{[-L,L]}e^{(\alpha-i\gamma)(t+L)}(t+L)^j,&\alpha<0,\\
 (2L)^{-1/2}1_{[-L,L]}e^{-i\gamma t}
                  ((t+L)/(2L))^j,&\alpha=0,
 \end{cases}
 \qquad 0\le j<m_\rho.
\]
Proposition~\ref{prop:hecke-deep-zero-jets} gives the squared
Hilbert--Schmidt bound $O(\Lambda^{4/3-56}(1+L)^C)$. Taking its
square root gives the assertion above, since every fixed negative
power of $\Lambda$ dominates every negative power of $L$.
\item For every fixed smooth strict-margin cutoff $p$, equal to one near
zero and supported in $[0,b]$ for some $b<1$, put
\[
 M_\Lambda=p(u/\Lambda),\qquad
 B_\Lambda^p=\mathcal H_cM_\Lambda\mathcal H_c,
 \qquad K_\Lambda^p=M_\Lambda B_\Lambda^pM_\Lambda.
\]
Proposition~\ref{prop:hankel-high-interior} gives, for every fixed
$N$, uniformly as an operator on the raw Hilbert space,
\begin{equation}\label{eq:hi-strict-margin-input}
 \|(1-\Pi_\Lambda)K_\Lambda^p\|_{L^2\to H_D^{3/8}}
                                      \le C_{p,N}\Lambda^{-N}.
\end{equation}
The two principal constraints can be restored exactly with the same
arbitrary-power graph accuracy. A concrete normalization of the correction
profiles and the elementary coordinate estimates are given below.
\item The Hankel transform preserves every $H_D^s$. On the critical line,
$|L_K(1/2+i\tau,\chi)|\ll(2+|\tau|)^{3/8}$, so
\begin{equation}\label{eq:hi-theta-graph}
       \|E_\chi f\|_2\le C_\chi\|f\|_{H_D^{3/8}}.
\end{equation}
The fixed-strip polynomial growth estimate, the Euler-product lower bound
at a fixed point to the right of one, and the zero count
$O_\chi(\log(2+|T|))$ per unit height are available.
\end{enumerate}
Only the strict-margin estimate, the fixed-packet annihilation, and the
analytic inputs are needed for the upper bounds. The rank, Gram, and
outside-window bounds explain their relation to the complete cutoff
construction and provide a trace-transfer corollary at the end.

\subsection{Support, scaling, and the principal constants}

For a nonprincipal character the raw time support gives
$E_\chi f(t)=0$ for $t>\log\Lambda$. The Poisson identity gives the
opposite support gap for a Hankel-band vector:
\[
 f\in\operatorname{Ran}B_\Lambda
       \quad\Longrightarrow\quad E_\chi f(t)=0\quad(t<-L).
\]
Indeed $\mathcal H_cf$ is time supported in $[0,\Lambda]$ and
$E_\chi f=\varepsilon_\chi J E_{\bar\chi}\mathcal H_cf$.
There is no conductor translation: the factor $c$ has already been
incorporated into the self-dual Hankel transform.

In the principal case let $r_K=\operatorname{Res}_{s=1}\zeta_K(s)$.
The theta multiplier normalization fixes the subtraction as
\[
 E_1f(t)=e^{t/2}\sum_{\mathfrak a\ne0}
                  f(N\mathfrak a\,e^t)-r_KIf\,e^{-t/2}.
\]
Here the ideal sum includes the unit ideal. The Euler factors give
$\zeta_K(s)=\zeta(s)L(s,\chi_{-3})$, and
\[
 L(1,\chi_{-3})
 =\int_0^1\frac{1-x}{1-x^3}\,dx
 =\int_0^1\frac{dx}{1+x+x^2}
 =\frac{\pi}{3\sqrt3}.
\]
Thus, for the primitive principal character, $c=2\pi/\sqrt3$ and
\begin{equation}\label{eq:hi-principal-residue}
                         r_K=\frac c6.
\end{equation}
The kernel at zero and involutivity give, on the smooth cores in use,
\begin{equation}\label{eq:hi-hankel-coordinates}
        (\mathcal H_cf)(0)=cIf,
        \qquad I(\mathcal H_cf)=\frac{f(0)}c.
\end{equation}
The principal root number is one. Consequently the two exact polar tails
are
\begin{align}
 E_1f(t)&=-r_KIf\,e^{-t/2} &&(t>L, P_\Lambda f=f),
 \label{eq:hi-upper-polar}\\
 E_1f(t)&=-\frac{r_K}{c}f(0)e^{t/2}
          =-\frac{f(0)}6e^{t/2}
                         &&(t<-L, B_\Lambda f=f).
 \label{eq:hi-lower-polar}
\end{align}
These support and polar identities extend to the graph domains by smooth
approximation preserving the indicated support and continuous coordinate.

The transform reverses raw scaling:
\[
 \mathcal H_cD_r=D_{1/r}\mathcal H_c,
 \qquad E_\chi D_r=\tau_{\log r}E_\chi,
 \quad (\tau_xg)(t)=g(t-x).
\]
Choose a fixed $\beta\in C_c^\infty((1/4,1/2))$ with $I\beta=1$,
and define the principal correction profiles
\begin{equation}\label{eq:hi-principal-profiles}
 b_{\Lambda,+}=D_\Lambda\beta,
 \qquad b_{\Lambda,-}=c^{-1}\mathcal H_cD_\Lambda\beta.
\end{equation}
For the normalized coordinates $I/\sqrt\Lambda$ and
$\operatorname{ev}_0/\sqrt\Lambda$, their coordinate matrix is exactly
the identity. Their theta images are fixed profiles after the appropriate
endpoint translation:
\[
 E_1b_{\Lambda,+}=\tau_LE_1\beta,
 \qquad E_1b_{\Lambda,-}=c^{-1}\tau_{-L}J E_1\beta.
\]
The factor $c^{-1}$ replaces the factor $1/2$ of the one-dimensional
even Fourier normalization.

\subsection{The Hankel cross-boundary estimate}

The only new localization calculation needed for the upper argument is
the replacement of ordinary Fourier tail estimates by Hankel tail
estimates. Set
\[
                  \mathcal D=-\partial_u(u\partial_u).
\]
The Bessel equation and integration by parts give
\begin{equation}\label{eq:hi-radial-diagonalization}
 \mathcal H_c\mathcal D f(v)=c^2v\mathcal H_cf(v),
 \qquad
 U\mathcal DU^{-1}=-e^{-t}(\partial_t-1/2)^2.
\end{equation}
Initially one may take $f\in C_c^\infty((0,\infty))$. The identity
extends whenever its displayed graph norms are finite, by truncation.
For the cutoff profiles flat near raw zero the boundary terms there also
vanish, so the same identity applies to them.

\begin{lemma}\label{lem:hi-hankel-cross}
Suppose a raw column operator $V$ has logarithmic output supported in
$t\ge-A$, and all fixed logarithmic output derivatives have finite
Hilbert--Schmidt norm. For every fixed $j,N$, its Hankel tail beyond a
spectral radius $Y$ satisfies
\[
 \|(1-p(v/Y))\mathcal H_cV\|_{2,H_D^j}
 \le C_{j,N,A}Y^{-N}
       \sum_{r\le j+2N}\|\partial_t^rUV\|_2.
\]
An enlargement of the finite derivative budget is harmless if $j$ is
not an integer. After translating a raw upper endpoint of radius
$\Lambda$ to radius one, the opposite cutoff has $Y=\Lambda^2$.
In particular, for the endpoint columns constructed below,
\begin{equation}\label{eq:hi-cross-errors}
 \|(1-B_\Lambda^p)(1-M_\Lambda)T\|_{2,H_D^{3/8}}
 +\|(1-M_\Lambda)(1-B_\Lambda^p)T\|_{2,H_D^{3/8}}
                         \le C_N\Lambda^{-N}.
\end{equation}
At height $H$, the right side becomes $H^{C_N}\Lambda^{-N}$ when
the endpoint derivative budgets are bounded by fixed powers of $H$.
\end{lemma}
\begin{proof}
For integer output derivative order, commute logarithmic differentiation
through the smooth spectral cutoff and use the reversal of the dilation
generator by $\mathcal H_c$. On the support of the spectral tail,
$v\gg Y$. Insert $v^{-N}$ and use
\eqref{eq:hi-radial-diagonalization} to transfer $v^N$ to
$c^{-2N}\mathcal D^N$ on the raw side. In logarithmic coordinates
$\mathcal D^N$ is $e^{-Nt}$ times a fixed differential polynomial of
order $2N$. The factors $e^{-Nt}$ are bounded on $t\ge-A$.
This proves the asserted estimate, column by column and then in
Hilbert--Schmidt norm. An integer derivative bound controls the required
fractional one.

By $\mathcal H_cD_\Lambda=D_{1/\Lambda}\mathcal H_c$, a cutoff at
the physical dual radius $\Lambda$ becomes a cutoff at
$Y=\Lambda^2$ after raw upper-endpoint scaling. This proves the first
term of \eqref{eq:hi-cross-errors}. Apply $\mathcal H_c$ to the second
term and use
\[
 \mathcal H_c(1-M_\Lambda)(1-B_\Lambda^p)
       =(1-B_\Lambda^p)(1-M_\Lambda)\mathcal H_c
\]
to reduce it to the same estimate for the Fourier-dual endpoint column.
\end{proof}

For later use, every fixed smooth raw profile supported strictly inside
$(0,1)$, dilated by $D_\Lambda$, is approximated in theta graph norm
by the selected image with arbitrary powers of $\Lambda^{-1}$.
Choose $p$ equal to one on its support. The preceding tail estimate
shows that $K_\Lambda^pD_\Lambda\phi-D_\Lambda\phi$ is arbitrarily
small in $H_D^{3/8}$; now use \eqref{eq:hi-strict-margin-input} and
\eqref{eq:hi-theta-graph}. The dual assertion follows by
$\mathcal H_c$ and the Poisson identity. It is enough to approximate
$\Pi_\Lambda K_\Lambda^pD_\Lambda\phi$; no estimate for
$\Pi_\Lambda$ on arbitrary fractional Sobolev vectors is required.

\paragraph{Restoring the principal coordinates.}
The profiles in \eqref{eq:hi-principal-profiles} admit the same graph
approximations. Their normalized coordinate matrix after selection tends
to the identity with arbitrary-power accuracy. Here point evaluation is
not inferred from $H_D^{3/8}$. For a selected vector $\psi$, its band
extension $w=B_\Lambda(P_\Lambda B_\Lambda P_\Lambda)^{-1}\psi$
satisfies $P_\Lambda w=\psi$ and $\|w\|\ll\|\psi\|$.
From the Bessel kernel and $|J_1(x)|\le x/2$,
\[
 |w(0)|\ll\Lambda^{1/2}\|\psi\|,
 \qquad \|w'\|_{L^2(0,1)}\ll\Lambda^{3/2}\|\psi\|.
\]
The smooth-cutoff columns have the same coarse derivative bound.
For their raw projection error $F$, use
\[
 |F(0)|^2\le C\|F\|_{L^2(0,1)}
                (\|F'\|_{L^2(0,1)}+\|F\|_{L^2(0,1)}).
\]
Cauchy--Schwarz over columns gives arbitrary-power smallness of the
normalized evaluation row. The mean row follows by Cauchy--Schwarz on
$[0,\Lambda]$; the exterior tail of the fixed dual profile is arbitrarily
small by \eqref{eq:hi-radial-diagonalization}. Subtracting the selected
profiles times the inverse of their two-by-two coordinate matrix restores
both constraints exactly. All constants depend on fixed derivative
budgets, not on the number of columns. This also records the uniformity
needed when the column norm grows with frequency height.

\subsection{Local equality}

For a nontrivial zero $\rho=\beta+i\gamma$ put
$\alpha=\beta-1/2$ and let $m_\rho$ be its multiplicity. Define
\begin{equation}\label{eq:hi-zero-measure}
 d\mathfrak b_\chi(\tau)=
 \sum_{\alpha=0}m_\rho\delta_\gamma(d\tau)
 +\sum_{\alpha\ne0}m_\rho
       \frac{|\alpha|\,d\tau}{\pi((\tau-\gamma)^2+\alpha^2)}.
\end{equation}
The unit-height zero count and $|\alpha|<1/2$ show that this is a
positive locally finite measure, with mass $O_\chi(\log(3+|T|))$ on
each unit interval. For distant ordinates, sum the inverse-square Cauchy
tails against the zero count. In particular it is integrable against
every Schwartz function.

\begin{lemma}[The one-sided Hecke closure]
\label{lem:hecke-halfline-closure}
The closure of the theta images of smooth raw profiles compactly
supported in $(0,1)$ is the inner multiple of the negative-half-line
Hardy space whose inner factor is the Blaschke product of the zeros
of $L_K(s,\chi)$ in $\Re s>1/2$, with their multiplicities.
For the principal character impose $I\psi=\psi(0)=0$.
The reflected Hankel-dual closure at the other endpoint has model
space given by the zeros in $\Re s<1/2$.
\end{lemma}
\begin{proof}
The theta--Poisson identity makes every smooth core image square
integrable, and raw support makes its logarithmic support nonpositive.
Raw contractions preserve the core and become left translations.
The closed image is therefore a nonzero translation-invariant subspace
of the half-line Hardy space. The invariant-subspace theorem writes it
as an inner multiple of that space, as in
Lemma~\ref{lem:local-halfline-closure}.

For a nonprincipal character the Mellin transforms of the generators
are $L_K(s,\chi)\mathcal M\phi(s)$. For the principal character take
$\psi=(u\partial_u+1)\phi$ with $\phi\in C_c^\infty((0,1))$.
Then $I\psi=\psi(0)=0$, and the transform is
$(1-s)\zeta_K(s)\mathcal M\phi(s)$. The factor at $s=1$ is regular
and has nonzero value $-r_K$. At any other prescribed point, a
sufficiently narrow bump has nonzero Mellin transform. Thus the common
zeros in the Hardy half-plane are exactly the nontrivial zeros, with
their orders.

Every generator continues analytically across each finite boundary
point. A common singular inner measure on that boundary is consequently
excluded: on a small interval where one generator is nonzero its inner
factor has no singular measure, while at its isolated boundary zeros
analytic finite-order vanishing also excludes a singular atom.
A remaining exponential inner factor would make every image vanish on
a fixed interval immediately below logarithmic zero. This is impossible.
Choose a bump supported arbitrarily close to raw $u=1$; for $e^t>1/2$
the theta sum has only the unit-ideal term. The same holds for the
principal differentiated bump, whose subtracted mean is zero.
There is no common delay, and the inner factor is precisely the stated
Blaschke product. Finally the functional equation and
$E_\chi\mathcal H_c=\varepsilon_\chi J E_{\bar\chi}$ identify the
dual right-half-plane divisor with the original left-half-plane divisor.
\end{proof}

\begin{theorem}\label{thm:hi-local-limit}
For the Hankel cutoff with the fixed exponent $M=56$, for every
$a\in C_c^\infty(\mathbb R)$,
\[
 \lim_{\Lambda\to\infty}\operatorname{Tr}(D_LC_aC_a^*)
                   =\int |a(\tau)|^2d\mathfrak b_\chi(\tau).
\]
The conclusion holds in the constrained principal sector, with the
constants in \eqref{eq:hi-principal-residue}--\eqref{eq:hi-lower-polar}.
No hypothesis on zero location, simplicity, or separation is imposed.
\end{theorem}
\begin{proof}
Write $F_\chi(\tau)=L_K(1/2+i\tau,\chi)$. Choose a fixed compact
interval containing $\operatorname{supp}a$ in its interior, with endpoints
away from critical zero ordinates. Let $P_{0,L}$ project onto the span of
the window jets of all critical-line zeros in that interval, and put
$W_L=S_L-P_{0,L}$.

Give every real pole of $b=a/F_\chi$ the negative-time prescription
$(\tau-\gamma-i0)^{-j}$. Its inverse transform $\kappa$ is smooth;
at positive infinity it and its derivatives are Schwartz, and at negative
infinity it is a finite exponential polynomial plus a Schwartz remainder.
This follows by subtracting the finitely many compactly cut off principal
parts: the inverse transform of a prescribed pole is a constant times
$1_{t<0}t^{j-1}e^{-i\gamma t}$.

Set $T_L=U^{-1}C_bW_L$. The input jet conditions annihilate every whole
exponential polynomial in its far negative output tail. The remaining
two exterior quadrant kernels are Hilbert--Schmidt, while the middle
kernel has polynomial growth in $L$. Thus for every fixed $j$,
\begin{equation}\label{eq:hi-fixed-bulk}
 \|UT_L\|_{\mathcal S_2(L^2,H^j)}\le C_j(1+L)^{d_j},
 \qquad E_\chi T_L=C_aW_L.
\end{equation}
The latter identity is exact in the closed multiplier realization.

The Mellin action of the Hankel transform is
\[
 \mathcal M(\mathcal H_cf)(s)
 =g_c(s)\mathcal Mf(1-s),
 \qquad g_c(s)=c^{1-2s}\frac{\Gamma(s)}{\Gamma(1-s)}.
\]
On the critical line the multiplier is smooth and unimodular.
Reflection changes the real-pole prescriptions; changing them back adds
only delta derivatives, which vanish on the reflected jet constraints.
Hence the dual column has the same kernel descriptions and polynomial
bounds. Equivalently one may use the exact Poisson identity and the dual
reciprocal multiplier.

After translating the upper endpoint to zero, the raw upper error has,
before its finite input-jet correction, kernel
\[
                 (1-p(e^t))\kappa(t-s)1_{[-2L,0]}(s).
\]
The output is bounded below by a fixed number and the input is
nonpositive. The positive Schwartz tail of $\kappa$ therefore gives a
fixed Hilbert--Schmidt limit, including every fixed logarithmic derivative
and output polynomial weight. The translated jet projections tend strongly
to zero. To check this directly, scale their normalized polynomial
amplitudes by $L^{-1/2}$; their fixed Gram matrices converge to nonsingular
Hilbert matrices, while distinct ordinates decouple by integration by
parts. Composition with the fixed Hilbert--Schmidt quadrant kernel then
removes their correction in Hilbert--Schmidt norm. The dual endpoint has
the identical property.

Lemma~\ref{lem:hi-hankel-cross} supplies the two cross errors with
arbitrary powers of $\Lambda^{-1}$. In a nonprincipal sector set
$A=M_\Lambda$, $B=B_\Lambda^p$, and $P_+=1_{t\le L}$. The exact
identity is
\begin{align*}
 S_LE_\chi T_L-E_\chi ABA T_L
   ={}&S_LE_\chi(1-A)T_L
      +S_LE_\chi(1-B)AT_L\\
     &+P_+E_\chi(1-A)BAT_L.
\end{align*}
Indeed expansion gives an extra $(P_+-S_L)E_\chi BAT_L$, which vanishes
by the band support gap, and the final time-supported image has no
component above $L$. The cross estimates replace $AT_L$ by $T_L$ in
the second term and $BAT_L$ by $AT_L$ in the last, with vanishing
Hilbert--Schmidt error after applying \eqref{eq:hi-theta-graph}.
The residual therefore has two fixed translated Hilbert--Schmidt limits.
Include also $S_LC_a(1-S_L)$, whose two outside-input quadrants have
fixed Hilbert--Schmidt limits because $\check a$ is Schwartz.

For the principal sector put $g=ABA T_L$,
$r_+=Ig/\sqrt\Lambda$, $r_-=g(0)/\sqrt\Lambda$, and
\[
 g^0=g-b_{\Lambda,+}r_+-b_{\Lambda,-}r_-.
\]
Its two coordinates vanish. Expansion of the preceding identity now has
the following exact form. If $R_0$ denotes its three-term right side,
\[
 R_0=S_LE_1T_L+(P_+-S_L)E_1BAT_L-P_+E_1g,
\]
so that
\[
 S_LE_1T_L-E_1g
 =R_0+r_KIg\,e^{-t/2}1_{t>L}
       +(r_K/c)g(0)e^{t/2}1_{t<-L}.
\]
The value of $BAT_L$ at zero equals that of $g$ because $A$ is
one there; the band-side polar tail is consequently expressed by
the same $g(0)$ as the constraint row.
Adding back the two correction profiles cancels these tails exactly,
leaving only
\begin{equation}\label{eq:hi-principal-boundaries}
 1_{t\le L}E_1b_{\Lambda,+}r_+
       +1_{t\ge-L}E_1b_{\Lambda,-}r_-.
\end{equation}

The two rows have fixed translated $L^2$ limits. For the mean row, move
$B-1$ to the normalized smooth integration test; its Hankel tail is
arbitrarily small. The main row is
$\Lambda^{-1/2}I(A^2T_L)$, whose translated kernel is the integral of
$\kappa(t-s)$ against $e^{t/2}p(e^t)^2$. For the value row use
$(\mathcal H_cf)(0)=cIf$ and move the first time cutoff to the smooth
Hankel integration test. The main dual row has weight
$c e^{t/2}p(e^t)$. These weights have exponential negative tails and
bounded upper support, so their row kernels have all fixed integrable
moments on $s\le0$. Their finite jet corrections vanish as above.
The output profiles in \eqref{eq:hi-principal-boundaries} are fixed after
translation and belong to $L^2$. Thus they are legitimate fixed
Hilbert--Schmidt boundary terms.

In both sectors we have constructed
\begin{equation}\label{eq:hi-endpoint-decomposition}
 S_LC_a(1-P_{0,L})
   =F_L+\tau_LB_+\tau_{-L}
          +\tau_{-L}B_-\tau_L+o_{\mathcal S_2}(1),
\end{equation}
where $B_+$ has range in $L^2(( -\infty,0])$, $B_-$ has range in
$L^2([0,\infty))$, and $\|F_L\|_2\le C(1+L)^d$.
The bulk is $E_\chi ABA T_L$ or $E_1g^0$, respectively.
The strict-margin input and exact constraint restoration give
\begin{equation}\label{eq:hi-bulk-distance}
                         \|(1-Q_\Lambda)F_L\|_2\longrightarrow0.
\end{equation}
This is a direct graph approximation by columns in the theta image;
it does not infer image-projection closeness from raw-projection closeness.

Let $M_+\subset L^2(( -\infty,0])$ be the one-sided theta closure and
$M_-\subset L^2([0,\infty))$ its dual reflected version.
Lemma~\ref{lem:hecke-halfline-closure} identifies their complements
$K_\pm$ as the pure Blaschke model spaces for the right and left
half-strip zeros of the original character.

The fixed-profile graph approximation proves
$\|(1-Q_\Lambda)\tau_{\pm L}v\|\to0$ for each $v\in M_\pm$.
Finite exponential-jet packets are dense in $K_\pm$. Their annihilation rate, and the exponentially small truncation of a fixed
off-line packet beyond the opposite endpoint, give for every fixed
packet projection $P_{K_{\pm,n}}$ and every $N$,
\[
 \|Q_\Lambda\tau_{\pm L}P_{K_{\pm,n}}\|
                                      =O_{n,N}(L^{-N}).
\]
Here $4/3-56=-164/3<0$ is more than sufficient; no height-uniform
packet Gram estimate is required.

For clarity, the endpoint Hilbert-space conclusion can be verified
directly from \eqref{eq:hi-endpoint-decomposition}. Strong convergence
on $M_\pm\oplus K_\pm$, followed by composition with the fixed
Hilbert--Schmidt boundaries, gives
\[
 (1-Q_\Lambda)S_LC_a(1-P_{0,L})
 =\tau_LP_{K_+}B_+\tau_{-L}
  +\tau_{-L}P_{K_-}B_-\tau_L+o_{\mathcal S_2}(1).
\]
Project the original decomposition onto a fixed endpoint packet.
Its bulk tends to zero since
\[
 \|P_{K_{+,n}}\tau_{-L}F_L\tau_L\|_2
 \le \|Q_\Lambda\tau_LP_{K_{+,n}}\|\,\|F_L\|_2
           +\|(1-Q_\Lambda)F_L\|_2\longrightarrow0.
\]
The opposite boundary tends to zero by compactness and weak convergence
of translations. The translated input projection tends strongly to zero.
Thus $P_{K_{+,n}}B_+=P_{K_{+,n}}C_a$, and exhaustion gives
$P_{K_+}B_+=P_{K_+}C_a$; the lower endpoint is identical.
The two translated boundaries are asymptotically Hilbert--Schmidt
orthogonal, by truncating their output ranges to fixed bounded intervals.

Lemma~\ref{lem:critical-window-packet} gives the critical contribution
$\sum_{\Re\rho=1/2}m_\rho|a(\Im\rho)|^2$ from the fixed-packet
annihilation. Its inverse-Gram calculation also accounts for repeated
critical zeros without commuting the multiplier and packet projection.
The critical input and its orthogonal complement add exactly in squared
Hilbert--Schmidt norm. Finally a Malmquist basis of each finite Blaschke
packet has squared Fourier diagonal equal to the sum
$2|\alpha|/((\tau-\gamma)^2+\alpha^2)$, including repeated zeros.
Increasing packets and the zero count give
\[
 \|P_{K_+}C_a\|_2^2+\|P_{K_-}C_a\|_2^2
 =\sum_{\alpha\ne0}m_\rho
       \int |a(\tau)|^2
             \frac{|\alpha|d\tau}{\pi((\tau-\gamma)^2+\alpha^2)}.
\]
The identity
$\|(1-Q_\Lambda)S_LC_a\|_2^2=\operatorname{Tr}(D_LC_aC_a^*)$
completes the proof.
\end{proof}

Polarization gives convergence on smooth compact frequency tests. A
smooth squared test equal to one on a fixed compact set bounds its masses;
uniform approximation then gives vague convergence of $\mu_\Lambda$ to
$\mathfrak b_\chi$. No noncompact test has yet been used.

\subsection{A reciprocal parametrix uniform in zero clusters}

Put $F_\chi(z)=L_K(1/2+iz,\chi)$ and $H=2+|T|$. A zero
$\rho=\beta+i\gamma$ corresponds to
$z_\rho=\gamma-i(\beta-1/2)$, so all nontrivial zeros lie in
$|\Im z|<1/2$.

\begin{lemma}\label{lem:hi-rectangle}
For sufficiently large $|T|$, there are a real bump
$a_T\in C_c^\infty((T-2,T+2))$, with $0\le a_T\le1$ and equal
to one on $[T-1/2,T+1/2]$, and a packet of $N\le C\log H$ zeros
$z_1,\ldots,z_N$, counted with multiplicity, with these properties.
Every zero projected into the support of $a_T$ belongs to the packet,
and $a_T$ is one near each such projection. If $R_{T,j}$ is the proper
rational sum of the principal parts of $z^j/F_\chi(z)$ at the packet
zeros, with negative-time prescriptions at real poles, then
\[
 r_{T,j}(\tau)=\frac{\tau^j a_T(\tau)}{F_\chi(\tau)}-R_{T,j}(\tau)
\]
is an ordinary smooth function and, for every fixed $j,m$,
\begin{equation}\label{eq:hi-remainder}
                  \|r_{T,j}^{(m)}\|_2\le H^{C_{j,m}}.
\end{equation}
If $R_{T,j,\pm}(t)$ are the inverse-transform exponential polynomials
on the positive and negative half-lines, extended as whole exponential
polynomials, then for each fixed $A$,
\begin{equation}\label{eq:hi-residue-difference}
 \sup_{-A\le t\le0}|R_{T,j,-}(t)-R_{T,j,+}(t)|
                                         \le H^{C_{j,A}}.
\end{equation}
The same assertions hold with an analytic numerator polynomially bounded
on a fixed neighborhood of the rectangle used in the proof. Constants
do not depend on distances between zeros or their distances to the axis.
\end{lemma}
\begin{proof}
Use a large fixed disk centered at $2+iT$ in the $s$-plane. The
fixed-strip bound is polynomial in $H$, the center value is bounded
below by the Euler product, and the number of zeros in a smaller disk
is $O(\log H)$. Divide by the disk Blaschke product of those zeros.
The quotient is zero free. The maximum principle bounds it above;
Harnack's inequality applied to a constant minus its logarithmic modulus
bounds its reciprocal on a fixed inner disk containing the required
rectangle in the $z$-coordinate. Fixed radii eight, between six and seven,
and five suffice. Large $|T|$ excludes trivial zeros and the principal
pole from these disks.

Let $\gamma_k$ be the real projections of the zeros in the Blaschke
disk, with multiplicity. On each of
$[T-3/2,T-1]$ and $[T+1,T+3/2]$, the integral of
$\log\min(1,|x-\gamma_k|)$ is bounded below uniformly. Summing and
averaging therefore gives $x_-,x_+$ with
\[
       \prod_k\min(1,|x_\pm-\gamma_k|)\ge e^{-CN_0}.
\]
Every factor is at most one, so each individual distance is at least
$e^{-CN_0}$. Choose a width $\varepsilon\ge H^{-C}$ small enough
that comparable product bounds hold within $4\varepsilon$ of the two
chosen lines. The vertical strips are zero free and have polynomial
reciprocal bounds. The horizontal sides at $\Im z=\pm3/4$ stay a
fixed distance from all nontrivial zeros. The resulting rectangle
$\Gamma$ has polynomial reciprocal bounds on its contour and suitable
side neighborhoods.

Take all zeros inside $\Gamma$ as the packet. Choose $a_T$ equal to
one between $x_-+\varepsilon$ and $x_+-\varepsilon$, and zero outside
$[x_--\varepsilon,x_++\varepsilon]$, with transitions in the zero-free
strips. Its fixed derivatives cost at most fixed powers of $H$.

Inside $\Gamma$, subtracting the principal parts leaves the Cauchy
integral of $z^j/F_\chi(z)$ on the contour. It and its fixed derivatives
are polynomially bounded away from a side. Within a transition strip,
move the vertical side to the other side of that strip without crossing
a zero. Cauchy's estimate, with distance comparable to $\varepsilon$,
gives the same bounds. This controls the real remainder and cancels every
real distributional pole, since the cutoff is one there.

For the tails, let $P(z)=\prod_{k=1}^N(z-z_k)$. Multiplying the contour
formula for the combined principal parts by $P$ gives a numerator of
degree less than $N$, with coefficients in powers of $z-T$ bounded by
$H^{C_j}C^N$. This follows by polynomial division on a fixed-size contour;
it involves no division by a root gap. Since $N=O(\log H)$, this is a
polynomial bound. Outside $|\tau-T|\le3$ the roots lie in a fixed
smaller rectangle, so
\[
 |R_{T,j}^{(m)}(\tau)|
       \le H^{C_{j,m}}(1+|\tau-T|)^{-m-1}.
\]
The remainder there is $-R_{T,j}$. This proves
\eqref{eq:hi-remainder}.

Closing the inverse-transform contour in the appropriate half-plane
shows
\[
 R_{T,j,-}(t)-R_{T,j,+}(t)
    =\frac1{2\pi}\int_\Gamma
                    \frac{z^j e^{-izt}}{F_\chi(z)}\,dz.
\]
On the fixed interval $-A\le t\le0$ the exponential is bounded by
$e^{3A/4}$, and the contour length is bounded. This proves
\eqref{eq:hi-residue-difference}. The analytic-numerator version uses
the same contour estimates.
\end{proof}

\begin{lemma}\label{lem:hi-allpole-boundaries}
Let $W_{L,T}$ be the orthogonal projection in $S_LL^2$ onto functions
whose entire Fourier transforms vanish at every packet zero $z_k$ to
its indicated multiplicity. Thus it removes the span of
$1_{[-L,L]}t^j e^{-i\overline z_k t}$. Put
\[
 b_T=a_T/F_\chi,
 \qquad T_{L,T}=U^{-1}C_{b_T}W_{L,T}.
\]
For every fixed $j$,
\begin{equation}\label{eq:hi-allpole-bulk}
 \|UT_{L,T}\|_{\mathcal S_2(L^2,H^j)}
             \le H^{C_j}(1+L)^{N+C_j},\qquad L\ge1.
\end{equation}
For a fixed smooth multiplier $m(t)$ supported in $t\ge-A$, with
bounded derivatives, the upper-translated boundary
$m(t)\tau_{-L}UT_{L,T}\tau_L$ has Hilbert--Schmidt norm at most
$H^C$, uniformly in $L$, after any fixed output derivative and polynomial
output weight. The dual endpoint has the same property.
\end{lemma}
\begin{proof}
The combined proper rational part admits a Newton expansion
\[
 R_{T,0}(z)=\sum_{j=1}^Nd_j\prod_{k=1}^j(z-z_k)^{-1},
 \qquad |d_j|\le H^C.
\]
The coefficient bound follows from the polynomial numerator in the
preceding proof, expressed in $z-T$; divided differences of a polynomial
use only bounded shifted roots and factors $C^N$, not root-gap
denominators.

On an input supported in $[-L,L]$ and with zero moment at $z_k$, divide
by $z-z_k$ using the decaying first-order exponential integration kernel.
Its modulus is at most one. The zero moment forces the result to vanish
beyond the opposite endpoint, so it remains supported in $[-L,L]$.
Its operator and Hilbert--Schmidt norms on that interval are at most
$2L$. Division removes one zero and preserves all remaining conditions,
including multiplicities. A product of $j$ inverse factors therefore
costs at most $(2L)^j$ in Hilbert--Schmidt norm. This also handles real
poles with the prescribed boundary value. The smooth remainder contributes
at most $(2L/(2\pi))^{1/2}\|r_{T,0}\|_2$. This proves
\eqref{eq:hi-allpole-bulk}; fixed derivatives cost only fixed powers of
$H$ because the output frequency support is $[T-2,T+2]$.

After upper-endpoint translation the input lies in $[-2L,0]$ and retains
all its exponential-jet moment conditions. For each output derivative,
use the decomposition with $R_{T,j}$ and $r_{T,j}$. The whole
exponential polynomial $R_{T,j,+}(t-s)$ annihilates the input projection
exactly. Subtract it first. The rational contribution left over is
\[
 1_{t-s<0}\bigl(R_{T,j,-}(t-s)-R_{T,j,+}(t-s)\bigr).
\]
On $t\ge-A$, $s\le0$, its support is the bounded triangle
$-A\le t\le s\le0$. Equation~\eqref{eq:hi-residue-difference}
therefore gives a polynomial Hilbert--Schmidt bound with every fixed
output weight. This is a bound on the combined residue polynomial,
even if individual residues become large.

For the smooth remainder, Plancherel and \eqref{eq:hi-remainder} give
all fixed weighted time norms. In particular
\[
 \int_{t\ge-A,\ s\le0}(1+|t|)^{2J}
              |\check r_{T,j}(t-s)|^2\,dt\,ds
 \le C_{J,A}\int_{u\ge-A}(1+|u|)^{2J+1}
              |\check r_{T,j}(u)|^2du\le H^{C_{J,j,A}}.
\]
Input restriction and right multiplication by an orthogonal projection
cannot increase the norm. Leibniz's rule handles $m$. Using the principal
part of $z^j/F_\chi$ at each derivative stage also includes correctly
any time-zero distribution produced by differentiation.

The exact Poisson identity identifies the dual inverse with the reflected
dual-character reciprocal. More explicitly,
$L_K(s,\bar\chi)g_c(s)=\varepsilon_{\bar\chi}L_K(1-s,\chi)$,
so
\[
        \frac{g_c(s)}{L_K(1-s,\chi)}
                    =\frac{\varepsilon_{\bar\chi}}{L_K(s,\bar\chi)}.
\]
Reflection sends the input zero conditions to exactly that dual packet.
Thus the dual boundary has the same estimates, with no gamma growth or
conductor displacement left over. Lemma~\ref{lem:hi-hankel-cross}
then gives the required ordinary radial derivative budgets as well.
\end{proof}

\subsection{Uniform local mass and Schwartz convergence}

\begin{theorem}\label{thm:hi-uniform-mass}
For the fixed exponent $M=56$ and the stated Hankel cutoff, there are
constants $C,A$ depending on the fixed character such that
\[
              \mu_\Lambda([T,T+1])\le C(2+|T|)^A
                     \qquad(\Lambda\ge2,\ T\in\mathbb R).
\]
The conclusion includes the constrained principal sector and requires no
zero-location or zero-separation assumption. Consequently
$\mu_\Lambda\to\mathfrak b_\chi$ against every Schwartz frequency
test.
\end{theorem}
\begin{proof}
Use the preceding bump centered at $T+1/2$. Let
$P_{L,T}=S_L-W_{L,T}$, which has rank at most $N\le C\log H$.
The local construction applies to $T_{L,T}$ with these additional
off-line zero conditions: the exact identity
$E_\chi T_{L,T}=C_{a_T}W_{L,T}$ is unchanged. By
Lemma~\ref{lem:hi-allpole-boundaries}, all boundary derivative and moment
budgets are bounded by fixed powers of $H$, uniformly in $L$. Their
theta images obey the same bounds by \eqref{eq:hi-theta-graph}.
The outside-input boundary $S_LC_{a_T}(1-S_L)$ is controlled by the
fixed derivative bounds of the bump.

In the principal sector the two constraint rows have the same polynomial
bound. Here are the details that avoid estimating individual residues.
For $k=1,2$, put $\theta_k(t)=e^{t/2}p(e^t)^k$. Then
$(\partial_t-1/2)\theta_k$ is smooth and compactly supported, and
\[
 \widehat\theta_k(z)
 =-\frac{\widehat{(\partial_t-1/2)\theta_k}(z)}{1/2+iz}.
\]
This continuation has only a possible pole at $z=i/2$. On the large-height
rectangles it is analytic with polynomial bounds. The normalized mean
row, up to an arbitrary-power cross error, is the restriction to $s\le0$
of
\[
       \check{(b_T\widehat\theta_2(-\cdot))}(-s)
\]
followed by the translated input projection. The whole reflected
positive-time residue polynomial is annihilated by that projection.
The analytic-numerator form of Lemma~\ref{lem:hi-rectangle} bounds the
remaining smooth part in $L^2$ by a fixed power of $H$. The value row
is the identical dual calculation with numerator
$c\widehat\theta_1(-\cdot)$. Only the unweighted $L^2$ norm of these
input rows is needed; orthogonal projection does not preserve uniform
polynomial input moments. The two output polar profiles in
\eqref{eq:hi-principal-boundaries} have fixed translated $L^2$ norms,
so the additional boundaries are bounded by $H^C$.

The strict-margin estimate is an operator estimate, with constants
independent of the column operator. Combining it with
\eqref{eq:hi-allpole-bulk}, and retaining the column norm explicitly,
bounds the bulk graph error after a fixed derivative budget by
\begin{equation}\label{eq:hi-uniform-bulk-error}
                 H^C(1+L)^{C\log H+C}e^{-\eta L}
\end{equation}
for a fixed $\eta>0$. Exact principal constraint restoration has the
same form by the coordinate estimates above. The cross errors are bounded
by this expression or by the stronger $H^Ce^{-\eta L}$.
Here the derivative budgets are fixed once; they do not grow with the
number of packet zeros. The fixed choice $M=56$ is unchanged.

Choose a fixed $D$ so large that all these errors are at most one when
$L\ge H^D$. Indeed the logarithm of \eqref{eq:hi-uniform-bulk-error}
is at most $C\log H+(C\log H+C)\log(1+L)-\eta L$, which is negative
uniformly above this threshold for sufficiently large $H$.
No approximation of the boundary ranges by the theta image is needed
for this upper bound; contractivity gives
\[
       \|(1-Q_\Lambda)S_LC_{a_T}(1-P_{L,T})\|_2\le H^C.
\]
The complementary input has squared norm at most
$\operatorname{rank}P_{L,T}\le C\log H$, since $\|a_T\|_\infty\le1$.
The input spaces are orthogonal, so their squared Hilbert--Schmidt norms
add exactly. Since $a_T=1$ on $[T,T+1]$, the distance identity proves the
desired mass bound for $L\ge H^D$.

For $L<H^D$, Bessel's inequality gives
\[
 0\le\frac{d\mu_\Lambda}{d\tau}\le\frac L\pi,
 \qquad \mu_\Lambda([T,T+1])\le H^D/\pi.
\]
Bounded heights follow from Theorem~\ref{thm:hi-local-limit} and the
same elementary estimate on the remaining bounded cutoff range.
This proves the uniform polynomial local mass bound.

For a Schwartz test, sum its rapidly decreasing supremum over unit
intervals against that polynomial mass bound. The tails are uniformly
summable. Vague convergence on the compact part therefore extends to
convergence against the full Schwartz test.
\end{proof}

\paragraph{Relation to the uncompressed trace.}
Theorem~\ref{thm:hecke-deep-leakage} also permits passage to the
uncompressed trace. If $Z_\Lambda$ is
the orthonormal theta column, then
$\|(1-S_L)Z_\Lambda\|_2=O(\Lambda^{-17/8})$ and its rank is
$O(\Lambda)$. Consequently
\[
 \|Q_\Lambda-S_LQ_\Lambda S_L\|_1
 \le2\sqrt{\operatorname{rank}Q_\Lambda}
                 \|(1-S_L)Z_\Lambda\|_2
                      =O(\Lambda^{-13/8}).
\]
Thus every bounded convolution operator has vanishing compression trace
error. For $h\in C_c^\infty(\mathbb R)$ and
$H(z)=\int h(t)e^{zt}dt$, the exact compressed identity and the preceding
Schwartz convergence give
\[
 2\log\Lambda\,h(0)-\operatorname{Tr}(Q_\Lambda V(h))
       \longrightarrow\int H(i\tau)\,d\mathfrak b_\chi(\tau).
\]
This proves the trace limit for the constrained principal image.
The next section supplies the two polar coordinates.
